\section{A single remaining family and two approaches to its complement}
\label{sec:frontier-update}

Fix the seven-step word of Theorem~\ref{thm:seven-three}. In this section
$B(t)$ denotes the physical state $(6t-1,1)$, or equivalently the
normalized state $(6t,1)$, with $t\ge1$. Write $F$ for a physical
update. Heights at a meeting are compared at the same phase. The
two clocks that reach that phase may differ by a multiple of seven.
The \href{https://github.com/bruscaj/perturbed_collatz/tree/main/certificates/frontier}{public proof supplements}
give the full recurrence and interval arguments, exact certificates and
instructions for reproducing each computation.

\subsection{Entry and induction on the original B1 parameter}
\label{sec:b1-complement}

\begin{theorem}\label{thm:b1-complement}
Every safe phased state is captured by the distinguished cycle if and only
if $B(t)$ is captured for each positive $t$ outside the five guards below.
\end{theorem}

To see why a single family suffices, write the normalized height of an
arbitrary target as $27v+r$, $1\le r\le27$. For each phase and each $r$ the
entry certificate supplies $A,B>0$ and $k\le30$ such that the launch with
parameter $Av+B$ reaches that exact target after $k$ updates for every
$v\ge0$. The checker substitutes the affine ray into each signed update,
checking its parity, phase, safety and endpoint. These $7\cdot27$ identities
cover every positive normalized height. An ancestor can be larger than
its target. Its relevance is that once its trajectory enters the cycle,
the target on that trajectory also reaches the cycle. The Lean proof uses
a smaller table: first advance to phase one in at most six updates, then
use one of the 27 phase-one inverse rays to that aligned state.

Define $J(t)$ by the following disjunction. All variables in these
conditions are nonnegative integers. Put
\[
\begin{split}
 E={}&\{(1,1),(2,3),(4,1),(5,3),(6,1),(8,3),(9,1),(11,3),(12,1),(13,3)\},\\
 R={}&\{9,11,12,13,25,27,28,31,36,45,47,51,52,55,\\
     &\hspace{1em}61,66,69,71,82,83,85,99,104,107,108,120,123,127\}.
\end{split}
\]
\begin{enumerate}
\item $3\mid t$, $t=2^kq$, and $(k\bmod14,q\bmod4)\in E$.
\item $3\mid t$, $t=2^{14h+6}q$, and $59049^hq\equiv15\pmod{32}$.
\item $t=2^{24}v+1$.
\item $t\bmod128\in R$.
\item For some $h,r>0$ and $q\ge0$,
\[
 601t+151=128^h q,\qquad 729^h q=601(1+131072\,128^h r)+151.
\]
\end{enumerate}

\begin{proof}
Assume capture outside $J$. Induct strongly on $t$. The first two guards
give a merger with $B(t/3)$, using respectively $k+3$ and $14h+12$
ticks on both clocks. Exact divisibility is essential here. The third
gives capture at $t=1$ in 18 ticks and otherwise a merger with
$B(1179648v+1)$ at clocks 25 and 18. The fourth gives a direct seven-tick
return to a smaller positive parameter. The fifth gives a merger with
$B(1+131072\,9^h r)$ at clocks $14h+18$ and 18. The strict comparison
with the initial $t$ follows from the two counter equations and is proved
in the repayment arithmetic module. These are precisely the progress
lemmas used by \texttt{global\_capture\_iff\_final\_obligation} in the
frontier Lean project. Determinism transfers capture across each merger.
The ancestor cover now gives capture for the full safe domain.
Conversely, every B1 launch outside $J$ is a safe state.
\end{proof}

The entire ray $t=128v+19$ lies outside $J$. This gives a concrete test
of the theorem's remaining hypothesis. It is a statement about this
five-guard predicate. For example, $t=2^{28}v+19$ on that very ray has
a smaller partner under the two-clock rule below. The older three-family
result classifies minimum counterexamples using $W$. Its reductions may
change phase and use a different family as the smaller launch. Induction
on $t$ within B1 requires a positive smaller B1 parameter.

\subsection{A legal exit, repayment, and a partial selector}
\label{sec:legal-exits}

The 88-row cover uses five phases. Its rank, in phase order
$(1,2,3,5,6)$, is
\[ Q=(4m,8m,6m,6m,12m+3). \]
Its five growing labels are
N00, N01, N02, S00, and N03. Its stopping rule uses phases $1,2,5$ for
launches in those phases and the full five-phase set for the other
launches. The separate first-return calculation using all five marks
throughout has 32 rows. The physical checker replays both partitions.
The 43-row three-family classifier uses another section and the rank $W$.

As a concrete $W$ reduction in A3, take $t=16u+4$ and suppose
\[
 115u+36=2^{7h+5}q,\qquad 3^h q\equiv7\pmod 8 .
\]
The original section state $(4t+3,3)$ and the alternate $(3t+2,5)$
meet after $7h+14$ and $7h+12$ ticks. Their initial rank difference is
$4(4t+3)-5(3t+2)=t+2$. The proof factors the paired trajectories
through a component of seven modes with one positive affine counter.
The congruence selects its terminal merger. At $h=0$, $u=84$ gives a
meeting, while $u=20$ has the required valuation but misses the terminal
residue. A bound on time spent inside this component leaves reentry and
the other exits to be accounted for.

A different construction composes B1 returns. If a positive parameter
admits $k\ge1$ legal growing macros followed by $\lceil k/4\rceil$
legal contracting macros of the certified G/H pair, the resulting
parameter is positive and strictly smaller. Each macro spans seven
physical ticks. The theorem \texttt{certified\_ceil\_repayment} checks
the resulting clock, safety and comparison against the original rank.
The public selector combines a wider set of arithmetic rules and
returns a capture certificate, a smaller-parameter certificate, or an
unhandled outcome. A separate recurrent pair grammar permits arbitrarily
many successful packets with an outer decreasing parameter. Its renewal
calculation assigns density about $19.254724\%$ to successful neighbor
mergers on $t=128v+19$ within that grammar. This percentage describes
its acceptance language on that progression. Universal availability of
successful certificates remains open.

\subsection{Coalescence for almost every neighbor and its obstruction}
\label{sec:pair-coalescence}

Garner and Gao studied heights and consecutive starts for the ordinary
map \cite{garner-heights,gao-heights}. The paired-recurrence supplement
also credits Duran and Leen for ordinary-map synchronization work and
Lyapunov Weaver for the stopped-prefix argument adapted to phased B1.
For a fixed arithmetic progression of $t\ge2$, the set of $t$ for which
$B(t)$ and $B(t-1)$ meet at the same clock has relative natural density
one. The written argument tracks a pair at common phase as
$(x,3^e x+b)$, ordering the two slopes. Each next parity of the free
coordinate selects one of two dyadic subrays with equal conditional
weight. On the resulting countable chain of pair types, an offset estimate
and an excursion estimate for the exponent give almost-sure return to a
finite set. Each member has a finite merging word with positive
conditional probability. Repeated visits therefore yield absorption
with probability one. Finite successful words give entire integer rays.
Finite unions of these rays have their cylinder densities, which increase
arbitrarily close to one. The proof of the drift, finite-set hitting and
the passage to arithmetic progressions is supplied in the paired
recurrence supplement. The proof uses a distribution on rays, without a
randomness hypothesis about an individual orbit.

A separate equal-grow-count selector leaves $O(X^{29/30})$ possible
minimum representatives through $X$. This bound counts parameters for
which that selector fails to find a smaller partner by its specified
clock. Even within the known basin, infinitely many parameters have
no smaller equal-count, equal-clock partner at any clock. For a
captured parameter the excess-grow index
$I(t)=\lim_{m\to\infty}(g_{7m}(B(t))-2m)$ stabilizes because the cycle
has two growing steps per lap. Explicit embeddings send $I$ to
$I-12j$ for arbitrary $j\ge0$. Taking the least captured parameter in
each of these distinct classes gives the obstruction, since an
equal-count synchronous merger preserves $I$.

Allowing unequal clocks yields, for every $v\ge0$, the following two
identities. Endpoints below are in normalized coordinates $(z,p)$.
\[
\begin{array}{c|c|c|c}
t&\text{smaller }u&\text{clocks from }B(t),B(u)&\text{endpoint}\\\hline
2^{39}v+17&16\cdot3^{20}v+2&(39,4)&(564859072962v+21,5)\\
2^{28}v+19&204073344v+16&(28,7)&(2324522934v+183,1)
\end{array}
\]
Direct signed replay on each entire affine ray checks both clocks and
all intermediate parities. The search is exhaustive within its two clock
budgets and its requirement that the partner stride be at most the source
stride. These examples illustrate how a sparse set of unresolved
representatives can require a different proof for its individual members.

\subsection{Ancestors in intervals and the boundary of the basin}
\label{sec:inverse-intervals}

\begin{theorem}\label{thm:inverse-interval}
For each fixed $\delta>0$ there is an effective $M_\delta$ such that every
phased target of normalized height $z>0$ has an ancestor $B(t)$ with
$X<t<(1+\delta)X$ whenever $X\ge M_\delta z$.
\end{theorem}
\begin{proof}
Align the target to phase one by at most six inverse shrinking steps.
Writing $a=91/127$, a seven-shrink lap multiplies $z-a$ by 128. For
each residue modulo three an available one-grow inverse lap multiplies
$z-a$ by $L=128/3$ with additive error of absolute value below 30.
A 28-tick connector with three grows then reaches a B1 state. Its 81
ternary residue cases are checked as whole physical rays. With $C$
equal to the aligned positive height minus $a$ and $K=2^{27}/81$,
$j$ shrinking laps and $g$ one-grow laps give
\[
 t=K L^g(128^j C+e),\qquad |e|<1.
\]
Indeed the normalized lap error is at most $30/(L-1)=18/25$ and the
connector contributes at most $1/10$. As $\log_{128}3$ is irrational,
finitely many $3^{-g}$ modulo powers of 128 form a multiplicative
net of any fixed fineness. For $0<\delta\le1$, choose maximum successive
ratio at most $1+\delta/4$, then place the leading term between
$(1+\delta/4)X$ and $(1+\delta/4)^2X$. Since $C\ge36/127$, choosing
$X/z$ sufficiently large ensures $j\ge0$ and relative error below
$\delta/16$. The resulting parameter is strictly inside the required
interval. For larger widths use $\delta=1$. The interval supplement
gives the explicit threshold and the finite connector.
\end{proof}

\begin{corollary}\label{cor:cone}
Global capture is equivalent to the existence of a fixed $\delta>0$
and an unbounded sequence of intervals $[X,(1+\delta)X]$ on each of
which all B1 parameters have the same capture label.
\end{corollary}
\begin{proof}
Global capture makes every label one. Otherwise apply
Theorem~\ref{thm:inverse-interval} to a captured target and an uncaptured
target. Each sufficiently large interval of that width contains B1
ancestors of both. They inherit different labels from their targets.
\end{proof}

The interval construction can be combined with a positive inverse-tree
potential at exponent $17/20$. A prefix-free cut in the inverse tree of
a nonperiodic target gives disjoint branches. Applying the interval
constructor below each cut leaf gives at least
$c_\delta(X/z)^{17/20}$ distinct parameters for sufficiently large $X/z$.
For a periodic target, first pass to an off-cycle ancestor of bounded
height. This count requires the precision-14 potential inequalities;
the single-ancestor theorem and the corollary do not. Both labels
would inherit these abundance estimates if an exceptional basin exists.
Proving that the known label occurs throughout intervals of a fixed
relative width along an unbounded sequence remains an open problem.


\subsection{An alternative criterion at phase two}
\label{sec:affine-root}
Let $P(n)$ mean capture of normalized height $n>0$ at phase two. There is
a second conditional equivalence: global capture holds exactly when
$P(n)\Leftrightarrow P(27n-22)$ for every $n\ge1$. The right-hand
argument is larger even at $n=1$. Equality here concerns capture labels;
it permits distinct trajectories with the same final label.

Seven ticks take the phase-two pair $(128v+17,128v+19)$ to
$(27v+5,v+1)$ at phase two. Substituting $n=v+1$ identifies the assumed
label equality with the missing adjacent-odd link. Under this hypothesis
the odd-pair certificate connects whole congruence
classes of odd neighbors at phase two. Cutting the ordered odds at the
remaining edges produces an initial singleton and 33 periodic block
types. Split by parameter parity. Of the 66 types, 37 admit a direct
descent, 19 split into 2,432 affine rays in total, and ten use chains
of previously certified equalities. Each successful block comparison
ends below the original block minimum or at a checked finite base.
Induction gives capture on the odd blocks. The signed path identities
and the reduction of the remaining phased states are detailed in the
affine-root supplement. Conversely, global capture makes both labels
true for every $n$. The proof uses written induction and exact affine
checks. Establishing the displayed label equality for all $n$ remains
a distinct open problem.


\subsection{Which subsets and which phase comparisons suffice?}
\label{sec:phase-balance}

The ancestor theorem also admits a restriction on the launch parameter.
Let $S$ be a set of positive integers recognized by a finite automaton in
a fixed integer base, with positive upper natural density. Every target
has a B1 ancestor with parameter in $S$. The regular-domain supplement
constructs synchronized blocks and places them by inverse steering in
digit intervals where the automaton guarantees an accepting member.
As an explicit example, take $S$ to be the integers whose binary digit
sum is even. The arithmetic checker covers all 137,781 phase and residue
cases; one of the selected launches reaches the target in at most 30
updates. Consequently, proving capture on this one digit language would
suffice. Entry is certified; capture on the language remains open.

There is also a criterion that compares the same normalized height
across phases. Write $P_i(z)\in\{0,1\}$ for the actual capture
indicator of $(z,i)$. For each fixed pair $i\ne j$, global capture is
equivalent to $P_i(z)=P_j(z)$ at every positive height. Under the
assumed equality, a proof may substitute one of those labels for the
other. The phase-quotient supplement then uses transport of equal
labels and, where necessary, a complete dyadic partition, finite bases
and induction on height. A substitution in this argument is a logical
step; it adds no available update to the original fixed dynamics.

For $i<j$ set $e_{ij}(x)=P_i(x+1)-P_j(x+1)$, $x\ge0$. Exact affine
paths give identities of the form
\[
 e_{ij}(x)+\sum_{\ell}e_{ij}(a_\ell x+b_\ell)=0,
 \qquad a_\ell\ge2,\quad 0\le b_\ell<a_\ell.
\]
Each identity telescopes signed differences after the two ends of
each common affine ray are advanced through their forced updates.
The checker verifies all 21 pairs, using 294 edges. A nonzero
difference at $x\ge1$ therefore forces an opposite sign at some
strictly larger argument. All phases at $x=0$ are captured. Repeating
the argument shows that eventual absence of either sign for one full
phase pair would suffice for global capture. The argument does not
establish that either sign is scarce.

For a more restricted comparison, define
$A(t)=\mathbf{1}\{P_2(3t+3)>P_6(3t+3)\}$. If
$P_1(18v+2)\ne P_1(18v+3)$ then one of 35 affine arguments gives
a positive error of this restricted kind. The two directed paths have
14 and 21 edges. Both directions matter since either adjacent label can
be the positive one. For integer $X\ge1$ this gives
$N_J(X)\le35N_A(8192X)$, where $N_J(X)$ counts those disagreements
for $0\le v<X$ and $N_A(X)=\sum_{0\le t<X}A(t)$. This pointwise
transfer supplies a counting inequality even while both supports are
unknown.


The extra implication from $J$ to global capture needs its own evidence.
A 63-state ternary machine gives three checked transitions per state and
literal capture paths for all 63 base heights. The correction terms on
each transition are sums modulo two of adjacent errors $J$. If $A(t)=0$
for all $t\ge0$, the two directed paths above force $J(v)=0$ for every
$v\ge0$. All transition corrections then vanish. Induction on ternary
digits gives capture at phase two, and every other phase advances to
phase two in finitely many steps. Conversely, global capture makes
$A$ identically zero. This yields the restricted one-sided criterion
with its exact progression and sign; universal absence of that error
remains open.
